\documentclass[conference]{IEEEtran}

\usepackage[utf8]{inputenc}
\usepackage[T1]{fontenc}
\usepackage{amsmath,amssymb}
\usepackage{booktabs}
\usepackage{tabularx}
\usepackage{cite}
\usepackage[hidelinks]{hyperref}

\newcommand{\SO}{\mathrm{SO}(3)}
\newcommand{\SE}{\mathrm{SE}(3)}
\newcommand{\avg}{\mathrm{avg}}

\begin{document}

\title{Generative Networks for Robot Grasping}

\author{
  \IEEEauthorblockN{Georgii Fatakhov}
  \IEEEauthorblockA{Technical University of Munich}
  \and
  \IEEEauthorblockN{Awar Satar}
  \IEEEauthorblockA{Technical University of Munich}
}

\maketitle

\begin{center}
  \small Proposal $\cdot$ ADLR WS 2026/27, project idea 4
\end{center}

\section{Topic and Main Idea}

Grasp synthesis for a parallel-jaw gripper is inherently multimodal, since a single object admits many valid
grasps. Generative models address this by learning a distribution over 6-DoF end-effector poses conditioned on the
observed object point cloud~\cite{eppner2021acronym,winkelbauer2022twostage}. Diffusion- and flow-based grasp
generators on $\SE$~\cite{urain2023se3dif,lim2024equigraspflow} achieve state-of-the-art quality, but they require
tens of network evaluations per sample, which limits their use in time-critical robotic applications.
MeanFlow~\cite{geng2025meanflow} and its extension to Riemannian manifolds~\cite{woo2026riemannianmeanflow} reduce
sampling to a few or a single network evaluation by learning the average velocity over a time interval instead of
the instantaneous velocity. Two recent works apply this principle to grasping. GraspMF~\cite{bukhari2026graspmf}
learns an average velocity field on $\SO \times \mathbb{R}^3$ and matches diffusion and flow baselines on ACRONYM
with at most five network evaluations, at up to 39 times lower latency. GraspMeanFlow~\cite{kwon2026graspmeanflow}
additionally preserves $\SE$-equivariance under few-step sampling and reaches with a single evaluation the
distribution coverage that an iterative $\SE$ flow model attains after five.

The goal of this project is to reimplement few-step, $\SE$-equivariant grasp generation, progressing from
EquiGraspFlow~\cite{lim2024equigraspflow} to GraspMeanFlow~\cite{kwon2026graspmeanflow}, on the ACRONYM
dataset~\cite{eppner2021acronym}, and to evaluate the resulting trade-off between grasp quality and inference
latency in simulation with MuJoCo. Semantically conditioned grasp generation and additional training data are
considered as optional extensions.

% Defined early so it floats to the top of page 2 (a table* never lands on the page it is defined on).
\begin{table*}[!t]
  \caption{Milestones}
  \label{tab:milestones}
  \centering
  \footnotesize
  \begin{tabularx}{\textwidth}{@{}l p{2.6cm} X X l@{}}
    \toprule
    \# & Milestone & Content & Completion criterion & Depends on \\
    \midrule
    M1 & EquiGraspFlow baseline & Codebase setup, ACRONYM preprocessing (4 categories), evaluation of the
      released checkpoint, own training & EMD of own training close to the published value and to the released
      checkpoint & -- \\
    M2 & MuJoCo evaluation & Floating parallel-jaw gripper (Panda hand), convex decomposition of object meshes,
      collision and contact checks, lift-and-hold test, latency measurement & Ground-truth grasps succeed at
      ${\gtrsim}\,90\,\%$, random grasps at ${\approx}\,0\,\%$; EquiGraspFlow evaluated across all NFE & -- \\
    M3 & GraspMeanFlow-SG & Two-time conditioning, boundary and semigroup loss, Euler sampler, NFE sweep &
      Equivariance verified to machine precision; EMD at NFE 1 clearly below that of EquiGraspFlow at NFE 1
      (reported: ${\approx}\,0.42$ vs.\ $0.73$); success rates evaluated in MuJoCo & M1, M2 \\
    M4 & Cloud scale-up & Transfer of training and evaluation to a GPU virtual machine on Google Compute Engine,
      with data and checkpoints in Cloud Storage; full NFE sweeps, ablations, and MuJoCo evaluation at scale &
      GMF-SG training on the virtual machine reproduces the local results; final experiments carried out there &
      M3 \\
    M5 & \textit{(optional)} Sampler improvement & Endpoint sampler with exponential $\SO$
      schedule~\cite{kwon2026graspmeanflow} & Higher success rate at NFE 5 than the Euler sampler & M3 \\
    M6 & \textit{(optional)} Extension & Part- or language-conditioned
      generation~\cite{lundell2023constrained}, or QDGset~\cite{huber2024qdgset} as additional training data &
      Defined upon selection & M3 \\
    \bottomrule
  \end{tabularx}
\end{table*}

\section{Technical Outline}

\textbf{Base model.} The base model is EquiGraspFlow~\cite{lim2024equigraspflow}, a conditional flow-matching model
on $\SO \times \mathbb{R}^3$ trained with Riemannian conditional flow
matching~\cite{lipman2023flowmatching,chen2024riemannianfm}. A VN-DGCNN encoder maps the object point cloud to an
equivariant feature, and a vector-neuron head predicts the angular and linear velocity of the grasp pose. The prior
is uniform on $\SO$ and Gaussian around the point-cloud centroid. Since the prior is $\SE$-invariant and the network
is $\SE$-equivariant by construction, generated grasps transform consistently with the object, and rotation
augmentation is not required in principle. We build on the public implementation and released checkpoint of
EquiGraspFlow.

\textbf{Few-step model.} For few-step generation, we reimplement the semigroup variant of
GraspMeanFlow~\cite{kwon2026graspmeanflow} (GMF-SG) on top of the EquiGraspFlow codebase, as no reference
implementation is publicly available. The network is extended to predict the average velocity
$(\omega^\avg, v^\avg)(s, t, T_t)$ over an interval $[s, t]$. The additional time input $h = t - s$ is injected
through the existing equivariant time-lifting mechanism, which preserves equivariance and adds a negligible number
of parameters. Training combines a flow-matching boundary term at $s = t$, which anchors the field to the data
distribution, with a semigroup consistency term, which requires a direct step over $[s, t]$ to agree with the
stop-gradient composition of two steps through an intermediate time $m \in [s, t]$. Both terms rely only on the
exponential and logarithm maps of $\SO$ and avoid Jacobian--vector products. At inference, a sample is obtained by
the update $R \leftarrow \exp(-\Delta t\,[\omega^\avg])\,R$ and $x \leftarrow x - \Delta t\, v^\avg$, with one
network evaluation per step; a single step yields one-step generation.

\textbf{Data.} We use ACRONYM~\cite{eppner2021acronym} with the four-category split of EquiGraspFlow (Bowl, Laptop,
Mug, Pencil). Test objects are unseen instances of the training categories and are presented in random
orientations.

\textbf{Evaluation.} Distribution coverage is measured by the earth mover's distance (EMD) between generated and
ground-truth grasp sets. As this metric is independent of the simulator, it permits direct comparison with
published results. Physical feasibility is first assessed geometrically by the collision rate, the fraction of
grasps in which the gripper intersects the object, and the contact rate, the fraction of grasps in which the object
lies between the fingers. It is then measured by the success rate of a lift-and-hold test in MuJoCo, which replaces
the Isaac Gym protocol of prior work. The test is calibrated with ground-truth grasps, which should succeed with
high probability, and with random grasps, which should fail.

Inference cost is reported as the number of network evaluations (NFE) and the wall-clock time per sample. Both
models are compared at matched $\text{NFE} \in \{1, 2, 5, 10, 20\}$ following the protocol
of~\cite{kwon2026graspmeanflow}, with EquiGraspFlow's default sampler as an additional reference point.

\textbf{Technical contribution.} The project contributes, first, an independent reimplementation of
GraspMeanFlow-SG~\cite{kwon2026graspmeanflow} from the publication alone, including a verification of the reported
few-step results. Second, it provides a MuJoCo-based grasp-success evaluation that is calibrated against
ground-truth and random grasps, applied identically to all methods, and independent of the deprecated Isaac Gym,
which also makes it a possible starting point for a later transfer to Isaac Sim. Third, it yields a controlled
comparison of flow matching and MeanFlow with identical backbone, data, and evaluation across the full range of
sampling budgets. Optional extensions are the endpoint sampler of~\cite{kwon2026graspmeanflow}, part- or
language-conditioned grasp generation related to constrained grasp sampling~\cite{lundell2023constrained}, and
additional training data from QDGset~\cite{huber2024qdgset}.

\section{Milestones}

Table~\ref{tab:milestones} lists the milestones. M1 and M2 are independent and are carried out in parallel. M3
constitutes the core result and requires both. Development begins locally, and M4 transfers the working pipeline to
Google Cloud for the larger experiments. Should M3 remain incomplete, M1 and M2 still yield an evaluated
flow-matching baseline.

\section{Alternative Topics}

As alternatives, we consider neural network-based world models (idea 11)~\cite{li2025robotworldmodel} and tactile
exploration of objects (idea 2)~\cite{watkinsvalls2019multimodal}.

\bibliographystyle{IEEEtran}
\bibliography{adlr}

\end{document}
